% from ega0-1-fr.tex
\oldpage[0\textsubscript{I}]{54}

\begin{env}[5.5.4]
\label{0.5.5.4-fr}
Soit $\mathcal{F}$ un $\mathcal{O}_X$-Module localement libre de rang $n$~;
il est immédiat que $\bigwedge^p\mathcal{F}$ est un $\mathcal{O}_X$-Module
localement libre de rang $\binom{n}{p}$ si $p\leq n$, réduit à $0$ si $p>n$,
car la question est locale et on est ramené au cas où
$\mathcal{F}=\mathcal{O}_X^n$~; en outre, pour tout $x\in X$,
\[
  (\bigwedge^p\mathcal{F})_x/
  \mathfrak{m}_x(\bigwedge^p\mathcal{F})_x
\]
est un espace vectoriel de dimension $\binom{n}{p}$ sur $k(x)$, qui
s'identifie canoniquement à
\[
  \bigwedge^p(\mathcal{F}_x/\mathfrak{m}_x\mathcal{F}_x).
\]
Soient
$s_1,\ldots,s_p$ des sections de $F$ au-dessus d'un ouvert $U$ de $X$, et
soit $s=s_1\wedge\cdots\wedge s_p$, qui est une section de
$\bigwedge^p\mathcal{F}$ au-dessus de $U$ (4.1.5)~; on a
$s(x)=s_1(x)\wedge\cdots\wedge s_p(x)$, et par suite, dire que
$s_1(x),\ldots,s_p(x)$ sont \emph{linéairement dépendants} signifie que
$s(x)=0$. On en conclut que \emph{l'ensemble des $x\in X$ tels que
$s_1(x),\ldots,s_p(x)$ soient linéairement indépendants est ouvert dans
$X$}~: il suffit en effet, en se ramenant au cas où
$\mathcal{F}=\mathcal{O}_X^n$, d'appliquer (5.5.2) à la section image de $s$
par une des projections de
$\bigwedge^p\mathcal{F}=\mathcal{O}_X^{\binom{n}{p}}$ sur les
$\binom{n}{p}$ facteurs.

En particulier, si $s_1,\ldots,s_n$ sont $n$ sections de $\mathcal{F}$
au-dessus de $U$ telles que $s_1(x),\ldots,s_n(x)$ soient linéairement
indépendantes en tout point $x\in U$, l'homomorphisme
$u:\mathcal{O}_X^n|U\to\mathcal{F}|U$ défini par les $s_i$ (5.1.1) est un
\emph{isomorphisme}~: on peut en effet se restreindre au cas où
$\mathcal{F}=\mathcal{O}_X^n$ et où on identifie canoniquement
$\bigwedge^n\mathcal{F}$ et $\mathcal{O}_X$~;
$s=s_1\wedge\cdots\wedge s_n$ est alors une section de $\mathcal{O}_X$
\emph{inversible} au-dessus de $U$, et on définit un homomorphisme réciproque
de $u$ au moyen des formules de Cramer.
\end{env}

\begin{env}[5.5.5]
\label{0.5.5.5-fr}
Soient $\mathcal{E}$, $\mathcal{F}$ deux $\mathcal{O}_X$-Modules localement
libres (de rang fini), et soit $u:\mathcal{E}\to\mathcal{F}$ un
homomorphisme. Pour qu'il existe un voisinage $u$ de $x\in X$ tel que $u|U$
soit \emph{injectif} et que $\mathcal{F}|U$ soit \emph{somme directe de
$u(\mathcal{E})|U$ et d'un sous-$(\mathcal{O}_X|U)$-Module localement libre
$\mathcal{G}$}, il faut et il suffit que
$u_x:\mathcal{E}_x\to\mathcal{F}_x$ donne, par passage aux quotients, un
homomorphisme \emph{injectif} d'espaces vectoriels
\[
  \mathcal{E}_x/\mathfrak{m}_x\mathcal{E}_x
  \longrightarrow
  \mathcal{F}_x/\mathfrak{m}_x\mathcal{F}_x.
\]
La condition est en effet \emph{nécessaire}, car $\mathcal{F}_x$ est alors
somme directe des $\mathcal{O}_x$-modules libres $u_x(\mathcal{E}_x)$ et
$\mathcal{G}_x$, donc $\mathcal{F}_x/\mathfrak{m}_x\mathcal{F}_x$ est somme
directe de
$u_x(\mathcal{E}_x)/\mathfrak{m}_x u_x(\mathcal{E}_x)$ et de
$\mathcal{G}_x/\mathfrak{m}_x\mathcal{G}_x$. La condition est
\emph{suffisante}, car on peut se borner au cas où
$\mathcal{E}=\mathcal{O}_X^m$~; soient $s_1,\ldots,s_m$ les images par $u$ des
sections $e_i$ de $\mathcal{O}_X^m$ telles que $(e_i)_y$ soit égal au $i$-ème
élément de la base canonique de $\mathcal{O}_y^m$ pour tout $y\in X$
(\emph{sections canoniques de $\mathcal{O}_X^m$})~; par hypothèse
$s_1(x),\ldots,s_m(x)$ sont linéairement indépendants, donc, si $\mathcal{F}$
est de rang $n$, il existe $n-m$ sections $s_{m+1},\ldots,s_n$ de
$\mathcal{F}$ au-dessus d'un voisinage $V$ de $x$ telles que les $s_i(x)$
$(1\leq i\leq n)$ forment une base de
$\mathcal{F}_x/\mathfrak{m}_x\mathcal{F}_x$. Il existe alors (5.5.4) un
voisinage $U\subset V$ de $x$ tel que les $s_i(y)$ $(1\leq i\leq n)$ forment
une base de $\mathcal{F}_y/\mathfrak{m}_y\mathcal{F}_y$ pour tout $y\in V$,
et on en conclut (5.5.4) qu'il y a un isomorphisme de $\mathcal{F}|U$ sur
$\mathcal{O}_X^n|U$, appliquant les $s_i|U$ $(1\leq i\leq m)$ sur les
$e_i|U$, ce qui achève la démonstration.
\end{env}

\section{Platitude}
\label{section:0.6-fr}

\begin{env}[6.0]
\label{0.6.0.0-fr}
La notion de platitude est due à J.-P.~Serre [16]~; dans ce qui suit, on omet
les démonstrations des résultats qui sont exposés dans l'\emph{Algèbre
commutative} de N.~Bourbaki, à laquelle nous renvoyons le lecteur. Nous
supposons tous les anneaux commutatifs.\footnote{Voir l'exposé cité de
N.~Bourbaki pour la généralisation de la plupart des résultats au cas non
commutatif.}

